\documentclass[10pt,a4paper]{amsart}
\usepackage{amsmath,amssymb,amsthm,graphicx}
\usepackage{verbatim}
\newenvironment{reference}{\comment}{\endcomment}
%\newenvironment{reference}{}{}
\newenvironment{slogan}{\comment}{\endcomment}
\newenvironment{history}{\comment}{\endcomment}

% For commutative diagrams we use Xy-pic
\usepackage[all]{xy}

% We use 2cell for 2-commutative diagrams.
\xyoption{2cell}
% Theorem environments.
%
\theoremstyle{plain}
\newtheorem{theorem}[subsection]{Theorem}
\newtheorem{proposition}[subsection]{Proposition}
\newtheorem{lemma}[subsection]{Lemma}

\theoremstyle{definition}
\newtheorem{definition}[subsection]{Definition}
\newtheorem{example}[subsection]{Example}
\newtheorem{exercise}[subsection]{Exercise}
\newtheorem{situation}[subsection]{Situation}

\theoremstyle{remark}
\newtheorem{remark}[subsection]{Remark}
\newtheorem{remarks}[subsection]{Remarks}

\numberwithin{equation}{subsection}

% Macros
%
\def\lim{\mathop{\mathrm{lim}}\nolimits}
\def\colim{\mathop{\mathrm{colim}}\nolimits}
\def\Spec{\mathop{\mathrm{Spec}}}
\def\Hom{\mathop{\mathrm{Hom}}\nolimits}
\def\Ext{\mathop{\mathrm{Ext}}\nolimits}
\def\SheafHom{\mathop{\mathcal{H}\!\mathit{om}}\nolimits}
\def\SheafExt{\mathop{\mathcal{E}\!\mathit{xt}}\nolimits}
\def\Sch{\mathit{Sch}}
\def\Mor{\mathop{\mathrm{Mor}}\nolimits}
\def\Ob{\mathop{\mathrm{Ob}}\nolimits}
\def\Sh{\mathop{\mathit{Sh}}\nolimits}
\def\NL{\mathop{N\!L}\nolimits}
\def\CH{\mathop{\mathrm{CH}}\nolimits}
\def\proetale{{pro\text{-}\acute{e}tale}}
\def\etale{{\acute{e}tale}}
\def\QCoh{\mathit{QCoh}}
\def\Ker{\mathop{\mathrm{Ker}}}
\def\Im{\mathop{\mathrm{Im}}}
\def\Coker{\mathop{\mathrm{Coker}}}
\def\Coim{\mathop{\mathrm{Coim}}}

% Boxtimes
%
\DeclareMathSymbol{\boxtimes}{\mathbin}{AMSa}{"02}

%
% Macros for moduli stacks/spaces
%
\def\QCohstack{\mathcal{QC}\!\mathit{oh}}
\def\Cohstack{\mathcal{C}\!\mathit{oh}}
\def\Spacesstack{\mathcal{S}\!\mathit{paces}}
\def\Quotfunctor{\mathrm{Quot}}
\def\Hilbfunctor{\mathrm{Hilb}}
\def\Curvesstack{\mathcal{C}\!\mathit{urves}}
\def\Polarizedstack{\mathcal{P}\!\mathit{olarized}}
\def\Complexesstack{\mathcal{C}\!\mathit{omplexes}}
% \Pic is the operator that assigns to X its picard group, usage \Pic(X)
% \Picardstack_{X/B} denotes the Picard stack of X over B
% \Picardfunctor_{X/B} denotes the Picard functor of X over B
\def\Pic{\mathop{\mathrm{Pic}}\nolimits}
\def\Picardstack{\mathcal{P}\!\mathit{ic}}
\def\Picardfunctor{\mathrm{Pic}}
\def\Deformationcategory{\mathcal{D}\!\mathit{ef}}
\usepackage{cleveref}
\crefname{originalreference}{original reference}{original references}
\Crefname{originalreference}{Original reference}{Original references}
\makeatletter
\newlabel{spaces-morphisms-definition-P}{{04RD}{source}{}{sourcepointer-1}{}}
\newlabel{spaces-morphisms-definition-P@cref}{{[originalreference][0][]04RD}{[1][0][]source}{}{}{}}
\makeatother
\title{Transferring \textbackslash{}'etale-local-on-source-and-target properties of morphisms of schemes to \textbackslash{}tau-local-on-the-source properties of morphisms of algebraic spaces}
\author{The Stacks Project Authors}
\date{stacks-project:spaces-descent.tex:local-snapshot-20261012; GNU FDL 1.2 or later}
\begin{document}
\maketitle
\noindent\textit{Editorial note.} All delivered mathematical text is the author's own, copied verbatim from the pinned Stacks Project chapter file spaces-descent.tex; no statement or proof step was rewritten, reordered or removed.\par\medskip
\noindent\textit{Editorial note.} The counted span cites one definition in another chapter, recorded in assembly.external\_references with canonical tags as navigation pointers.\par\medskip
\noindent\textit{Allowed standard prerequisites (editorial).}\begin{itemize}
\item The property P\_spaces attached to a scheme property. Stacks Project, Morphisms of Spaces, Definition spaces-morphisms-definition-P. Supplies the outer equivalences of the displayed chain.
\item Base changes and refinements of \textbackslash{}tau-coverings. Stacks Project, Topologies on Spaces, section on the \textbackslash{}tau-topologies. Underlies the covering bookkeeping in the middle step.
\end{itemize}\medskip
\ifdefined\hypertarget\hypertarget{sourcepointer-1}{}\fi
\noindent\textit{Source reference.} Original reference 04RD: sources/stacks-chow/tags.txt, source lines 6181--6181. Cross-chapter standard prerequisite cited by the author inside the counted span; it is used as a pointer and its own proof is not part of this item.\par\medskip
\begin{lemma}
\label{lemma-transfer-from-schemes}
Let $S$ be a scheme.
Let $\tau \in \{fpqc, \linebreak[0] fppf, \linebreak[0] syntomic, \linebreak[0]
smooth, \linebreak[0] \etale\}$.
Suppose that $\mathcal{P}$ is a property of morphisms of schemes over $S$
which is \'etale local on the source-and-target. Denote $\mathcal{P}_{spaces}$
the corresponding property of morphisms of algebraic spaces over $S$, see
Morphisms of Spaces, Definition \ref{spaces-morphisms-definition-P}.
If $\mathcal{P}$ is local on the source for the $\tau$-topology, then
$\mathcal{P}_{spaces}$ is local on the source for the $\tau$-topology.
\end{lemma}

\begin{proof}
Let $f : X \to Y$ be a morphism of algebraic spaces over $S$.
Let $\{X_i \to X\}_{i \in I}$ be a $\tau$-covering of algebraic spaces.
Choose a scheme $V$ and a surjective \'etale morphism $V \to Y$.
Choose a scheme $U$ and a surjective \'etale morphism $U \to X \times_Y V$.
For each $i$ choose a scheme $U_i$ and a surjective \'etale morphism
$U_i \to X_i \times_X U$.

\medskip\noindent
Note that $\{X_i \times_X U \to U\}_{i \in I}$ is a $\tau$-covering.
Note that each $\{U_i \to X_i \times_X U\}$ is an \'etale covering,
hence a $\tau$-covering. Hence $\{U_i \to U\}_{i \in I}$ is a
$\tau$-covering of algebraic spaces over $S$. But since $U$ and each $U_i$
is a scheme we see that $\{U_i \to U\}_{i \in I}$ is a
$\tau$-covering of schemes over $S$.

\medskip\noindent
Now we have
\begin{align*}
f \text{ has }\mathcal{P}_{spaces}
& \Leftrightarrow
U \to V \text{ has }\mathcal{P} \\
& \Leftrightarrow
\text{each }U_i \to V \text{ has }\mathcal{P} \\
& \Leftrightarrow
\text{each }X_i \to Y\text{ has }\mathcal{P}_{spaces}.
\end{align*}
the first and last equivalence by the definition of
$\mathcal{P}_{spaces}$ the middle equivalence because we assumed
$\mathcal{P}$ is local on the source in the $\tau$-topology.
\end{proof}

\end{document}
